\documentclass{article}
\usepackage[utf8]{inputenc}
\usepackage[T1]{fontenc}
\usepackage{geometry}
\geometry{
    a4paper,
    left=2.5cm,
    right=2.5cm,
    top=2.5cm,
    bottom=2.5cm
}
\usepackage{graphicx}
\usepackage{float}
\usepackage{hyperref}
\usepackage[english]{babel}
\usepackage{amsmath}
\usepackage{amssymb}
\usepackage{array}
\usepackage{booktabs}
\usepackage{caption}
\usepackage{subcaption}
\usepackage{enumitem}
\usepackage{fancyhdr}
\usepackage{titlesec}
\usepackage{url}
\usepackage{textcomp}
\usepackage{xcolor}
\usepackage{listings}
\usepackage{multirow}
\usepackage{longtable}
\usepackage{tabularx}
\usepackage{times}

% Set up page style
\pagestyle{fancy}
\fancyhf{}
\rfoot{\thepage}
\renewcommand{\headrulewidth}{0pt}

% Title formatting
\titleformat{\chapter}[display]{\normalfont\bfseries}{\chaptertitle\thechapter}{0pt}{\Huge}
\titlespacing*{\chapter}{0pt}{50pt}{40pt}

% Document info
\title{Template Document}
\author{Academic Support Agent}
\date{\today}

\begin{document}
\maketitle

%% The abstract of the dissertation in English
% Use this command!
\draftabstract{\lipsum[1-3]}%
% Let's add another one in Finnish
\draftabstract[finnish]{\hspace{-2pt}T ss  ty ss  positroniannihilaatiospektroskopiaa k ytettiin periodisten rakenteiden 
pistevirheiden tutkimiseen merkitt viss  ydinteknillisiss  materiaaleissa. T m n 
lis ksi ilmaisinten vaikutusta tuloksiin tutkittiin sek  elinaika- ett  
Doppler-levenem \-spektroskopiassa.

Ilmaisimilla on merkitt v  vaikutus elinaika- ja Doppler-levenem \-spektroskopioiden
suorituskykyyn ja saatavien tulosten laatuun. Elinaikaspektroskopia on hyvin herkk 
mitattavien n ytteiden ylim r iselle s teilylle, mik  rajoittaa sen k ytt  
ydinteknillisten materiaalien tutkimuksessa. Ainoa tapa lis t  mittauslaitteiston 
s teilytoleranssia on ottaa k ytt n kolmas koinsidenssi-ilmaisin. Doppler-levenem laitteistossa
s teilyilmaisimen energiaresoluutio vaikuttaa suoraan mittaustuloksiin.

Kahden ilmaisimen elinaikaspektroskopian herkkyytt  v rille koinsidenssitapahtumille 
tutkittiin lis m ll  koboltti-60 -l hde tutkittavien piin ytteiden viereen. 
S teilytoleranssin kasvua kolmen ilmaisimen elinaikaspektroskopiassa arvioitiin 
yksinkertaisilla teoreettisilla malleilla. Doppler-levenem \-spektroskopiassa tutkittiin 
energiaresoluutiota liitt m ll  natrium-22 -positronil hde suoraan tutkittaviin n ytteisiin,
jolloin pystyttiin samanaikaisesti mittaamaan sek  Doppler-leven m spektri  ett  
yksienergist  fotoabsorptiospektri . T ss  ty ss  tutkittiin l hde-ilmaisin -geometrian
vaikutuksia tuloksiin ja keskustellaan muuttuvien tulosten syist  perustuen systemaattisiin 
kokeisiin ja Monte Carlo -simulaatioihin.

Mikroskooppisilla virheill  on merkitt v  vaikutus ydinmateriaalien ominaisuuksiin, 
kuten korroosioon ja s teilykest vyyteen. Zircaloy-4 on yksi eniten k ytetyist  
polttoaineen suojakuoriseoksista ydinvoimalaitoksissa, mutta sen hapettuminen
reaktoriolosuhteissa on monimutkainen prosessi, johon kuuluu hapettumisnopeuden 
periodisuus. Fuusiovoimalaitoksissa plasmaa l hell  olevat materiaalit joutuvat
 rimm isen ankariin olosuhteisiin, joissa korostuvat korkea l mp tila ja suuret 
hiukkass teilyannokset. Volframia pidet n lupaavana materiaalikandidaattina t h n
tarkoitukseen.

Zirkoniumoksidin perustavanlaatuisia ominaisuuksia tarkasteltiin perustuen 
mikroskooppisten hilavirheiden muuttumiseen tutkittavissa Zircaloy-4 -n yt\-teiss , joita
oli hapetettu painevesireaktoria vastaavissa olosuhteissa. Doppler-levenem \-spektroskopiaa
ja teoreettista mallinnusta sovellettiin hapettumiskerrosten hilavirheiden tutkimiseen. 
Positronien elinaikaspektrokopiaa k ytettiin toivutettujen ja protonis teilytykselle 
altistettujen volframin ytteiden monovakanssien tutkimiseen. V lisija-atomien
ja monovakanssien migraatioenergiarajat pystyttiin havaitsemaan suoraan 
elinaikaspektrometrilla, joka oli yhdistetty kylm s teilytyslaitteistoon.

}%


% And yet another one in Swedish
\draftabstract[swedish]{\lipsum[7-9]}
%%---------------------

%% The abstract of the dissertation in English
% Use this command!
\begin{abstract}\lipsum[1-3]\end{abstract}%
% Let's add another one in Finnish
\begin{abstract}[finnish]T ss  ty ss  positroni\-annihilaatio\-spektroskopiaa k ytettiin periodisten rakenteiden 
piste\-virheiden tutkimiseen merkitt viss  ydinteknillisiss  materiaaleissa. Lis ksi tutkittiin ilmaisinten 
vaikutusta tuloksiin sek  elinaika- ett  Doppler-levenem \-spektro\-skopiassa.

Ilmaisimilla on merkitt v  vaikutus elinaika- ja Doppler-levenem \-spektroskopioiden
suorituskykyyn ja saatavien tulosten laatuun. Elinaikaspektroskopia on hyvin herkk 
mitattavien n ytteiden ylim r iselle s teilylle, mik  rajoittaa sen k ytt  
ydinteknillisten materiaalien tutkimuksessa. Ainoa tapa lis t  mittauslaitteiston 
s teilytoleranssia on ottaa k ytt n kolmas koinsidenssi-ilmaisin. Doppler-levenem laitteistossa
s teilyilmaisimen energiaresoluutio vaikuttaa suoraan mittaustuloksiin.

Kahden ilmaisimen elinaikaspektroskopian herkkyytt  v rille koinsidenssitapahtumille 
tutkittiin lis m ll  koboltti-60 -l hde tutkittavien piin ytteiden viereen. 
S teilytoleranssin kasvua kolmen ilmaisimen elinaikaspektroskopiassa arvioitiin 
yksinkertaisilla teoreettisilla malleilla. Doppler-levenem \-spektroskopiassa tutkittiin 
energiaresoluutiota liitt m ll  natrium-22 -positronil hde suoraan tutkittaviin n ytteisiin,
jolloin pystyttiin samanaikaisesti mittaamaan sek  Doppler-leven m spektri  ett  
yksienergist  fotoabsorptiospektri . T ss  ty ss  tutkittiin l hde-ilmaisin -geometrian
vaikutuksia tuloksiin ja keskustellaan muuttuvien tulosten syist  perustuen systemaattisiin 
kokeisiin ja Monte Carlo -simulaatioihin.

Mikroskooppisilla virheill  on merkitt v  vaikutus ydinmateriaalien ominaisuuksiin, 
kuten korroosioon ja s teilykest vyyteen. Zircaloy-4 on yksi eniten k ytetyist  
polttoaineen suojakuoriseoksista ydinvoimalaitoksissa, mutta sen hapettuminen
reaktoriolosuhteissa on monimutkainen prosessi, johon kuuluu hapettumisnopeuden 
periodisuus. Fuusiovoimalaitoksissa plasmaa l hell  olevat materiaalit joutuvat
 rimm isen ankariin olosuhteisiin, joissa korostuvat korkea l mp tila ja suuret 
hiukkass teilyannokset. Volframia pidet n lupaavana materiaalikandidaattina t h n
tarkoitukseen.

Zirkoniumoksidin perustavanlaatuisia ominaisuuksia tarkasteltiin perustuen 
mikroskooppisten hilavirheiden muuttumiseen tutkittavissa Zircaloy-4 -n yt\-teiss , joita
oli hapetettu painevesireaktoria vastaavissa olosuhteissa. Doppler-levenem \-spektroskopiaa
ja teoreettista mallinnusta sovellettiin hapettumiskerrosten hilavirheiden tutkimiseen. 
Positronien elinaikaspektrokopiaa k ytettiin toivutettujen ja protonis teilytykselle 
altistettujen volframin ytteiden monovakanssien tutkimiseen. V lisija-atomien
ja monovakanssien migraatioenergiarajat pystyttiin havaitsemaan suoraan 
elinaikaspektrometrilla, joka oli yhdistetty kylm s teilytyslaitteistoon.

\end{abstract}


% And yet another one in Swedish
\begin{abstract}[swedish]\lipsum[7-9]\end{abstract}


%% Preface
% If you write this somewhere else than in Helsinki, use the optional location.
\begin{preface}[Optional location (if not defined, Helsinki)]
\lipsum[1-4]
\end{preface}

%% Table of contents of the dissertation



% To be defined before generating list of publications. Leave off if no acknowledgement
\languagecheck{the Institute of Language Checks}

%% This is for article dissertations. Remove if you write a monograph dissertation.
% The actual publications are entered manually one by one as shown further down:
% use \addpublication, \addcontribution, \adderrata, and addpublicationpdf.
% The last adds the actual article, the other three enter related information
% that will be collected in lists -- like this one.
\listofpublications

%% Add lists of figures and tables as you usually do (



%% Add list of abbreviations, list of symbols, etc., using your preferred package/method.

\abbreviations

\begin{description}
\item[GPS] Global Positioning System
\item[ESA] European Space Agency
\item[WGS84] World Geodetic System 1984
\item[KKJ] Kartastokoordinaattij rjestelm 
\end{description}

\symbols

\begin{description}
\item[$W$] geopotential
\item[$G$] Newton's universal gravitational constant
\item[$c$] speed of light
\item[$h$] Planck's constant
\end{description}


%% The main matter, one can obviously use \input or \include
\chapter{Chapter Heading}
\section{Section Heading}
\[
\vec{F}=-G m_1 m_2 \frac{\vec{x}_2-\vec{x}_1}{\lVert{\vec{x}_2-\vec{x}_1}\rVert^3}
\]
\lipsum[1-2]
Kirjallisuusviite: \cite{Knuth1984:The-TeXbook}

Lause jossa on \textsf{sans-serif} ja \texttt{kirjoituskone} sanoja.

J kauden aikana noin 3 km paksu j vuori peitti alleen koko Skandinavian ja
osan Keski-Eurooppaa. J  painoi maanpintaa alas. Kun j  alkoi sulaa noin 12
000 vuotta sitten, maa nousi takaisin, ja nousu jatkuu yh . Se on 30--100 cm
sadassa vuodessa. Suurimmillaan se on Merenkurkussa; Pohjois-Pohjanmaan
rannikolla se on 88 cm, sis maassa v hemm n, mist  syyst  jokien lasku mereen
on vaikeutunut ja maa soistunut. T m  ilmi  on ongelma, joka jatkuu edelleen.
Toisaalta sen ansiosta Pohjois-Pohjanmaan maa-alue kasvaa muita maakuntia
nopeammin. Maannoususta runoilija Sakari Topelius sanoi 1800-luvulla, ett 
ihmiset saavat Pohjanmaalla uutta maata ruhtinaskunnan verran sadassa vuodessa.

\begin{table}
\caption{A commutative diagram.}
\begin{tabular}{ccc}
\hline
Space domain & Fourier & Frequency domain\tabularnewline
\hline
$V\left(x,y,0\right)$ & $\mathcal{F}\Longrightarrow$ & $v_{jk}$\tabularnewline
$\downarrow$ (hard) &  & $\Downarrow\times$ (easy)\tabularnewline
$V\left(x,y,z\right)$ & $\Longleftarrow\mathcal{F}^{-1}$ & $v_{jk}\exp\left(-\pi\sqrt{j^{2}+k^{2}}\frac{z}{L}\right)$\tabularnewline
\hline
\end{tabular}
\end{table}

% Useful package
\begin{quoting} 
T ll , jossa me nyt asumme, on siis muinoin ollut pauhaava
meri, ja viel kin l ytyy monin paikoin meriraakun kuoria maasta. 

Mutta se oli Jumalan tahto, ett  t m n maan piti kohiaman yl s merest  niin
ihmeellisell  tavalla, kuin harvat muut maat, paitsi Pohjois-Ruotsia.
\end{quoting}
What an extraordinary situation is that of us mortals! Each of us is here for a
brief sojourn; for what purpose he knows not, though he sometimes thinks he
feels it. But from the point of view of daily life, without going deeper, we
exist for our fellow-men~in the first place for those on whose smiles and
welfare all our happiness depends, and next for all those unknown to us
personally with whose destinies we are bound up by the tie of sympathy. A
hundred times every day I remind myself that my inner and outer life depend on
the labours of other men, living and dead, and that I must exert myself in
order to give in the same measure as I have received and am still receiving. I
am strongly drawn to the simple life and am often oppressed by the feeling that
I am engrossing an unnecessary amount of the labour of my fellow-men. I regard
class differences as contrary to justice and, in the last resort, based on
force. I also consider that plain living is good for everybody, physically and
mentally. 

\begin{figure}

***************

* F I G U R E *

***************

\caption{Example of a figure.}
\end{figure}

\section{Section Heading}
$E=mc^2$
\lipsum[5-6]
% Content from dummymathcode/testcode.tex would be included here

\cite{WikiBooks2008:LaTeX}
\subsection{Subsection Heading}
$\Delta x\Delta p \gtrapprox \hbar$
\lipsum[7-8]
\subsection{Subsection Heading}
\lipsum[9-10]
\section{Section Heading}
\lipsum[11-12]

%% Examples of article references. Replace by your own as appropriate!

 % Refer to the Journal paper 1 of this example document
\citepub{j1} \& \cpub{j1} \& \cp{j1} \& \pageref{j1} \& \ref{j1}

% Refer to the Conference paper of this example document
\citepub[p.~2]{c1} \& \cpub[Sec.~ 1]{c1} \&  \cp[pp.~1--2]{c1} \& \pageref{c1} \& \ref{c1} 


\chapter{Chapter Heading 2}
\section{Sample Mathematics}

% Content from dummymathcode/testmath.tex would be included here


\section{Section Heading}
\lipsum[1-4]
\section{Section Heading}
\lipsum[5-6]
\subsection{Subsection Heading}
\lipsum[7-8]
\subsection{Subsection Heading}
\lipsum[9-10]
\section{Section Heading}
\lipsum[11-12]

%% An example for changing the running header (the optional parameter)
\chapter[Short Chapter Heading]{Chapter Heading 3}
\section{Section Heading}
\lipsum[1-4]
\section{Section Heading}
\lipsum[5-6]
\subsection{Subsection Heading}
\lipsum[7-8]
\subsection{Subsection Heading}
\lipsum[9-10]
\section{Section Heading}
\lipsum[11-22]

\renewcommand{\bibname}{References}
 % Change as required
\LARGE  % remember to edit the file name


%% The following commands are for article dissertations, remove them if you write a monograph dissertation.

% Errata list, if you have errors in the publications.
\errata

%% The first publication (journal article)
% Set the publication information.
% This command musts to be the first!
\addpublication{Journal Paper Authors}{Journal Paper Title}{Journal Name}{Volume, issue, pages, and other detailed information}{Month}{Year}{Copyright Holder}{j1}
% Add the dissertation author's contribution to that publication (the order can be interchanged with \adderrata).
\addcontribution{The author did this and that}
% Add the errata of the publication, remove if there are none (the order can be interchanged with \addauthorscontribution).
\adderrata{j1 I This is wrong}
% Add the publication pdf file, the filename is the parameter (must be the last).
\addpublicationpdf{dummyarticles/dummypdfarticle1.pdf}

%% The second publication (conference article, note the optional parameter)
% Set the publication information.
\addpublication[conference]{Conference Paper Authors}{Conference Paper Title}{Conference Name}{Location, pages, and other detailed information}{Month}{Year}{Copyright Holder}{c1}
% Add the dissertation author's contribution to that publication.
\addcontribution{The author did also this and that}
% No errata
% Add the publication pdf file, the filename is the parameter.
\addpublicationpdf{dummyarticles/dummypdfarticle2.pdf}

%% The third publication (another journal paper, accepted for publication, note the optional parameter)
% Set the publication information, detailed information can be empty
\addpublication[accepted]{Journal Paper 2 Authors}{Journal Paper 2 Title}{Journal Name 2}{}{Month}{Year}{Copyright Holder}{j2}
% Add the dissertation author's contribution to that publication.
\addcontribution{The author did everything}
% Add the errata of the publication, remove if there are none.
\adderrata{j2 III This is wrong}
% Add the publication pdf file, the filename is the parameter.
\addpublicationpdf{dummyarticles/dummypdfarticle3.pdf}

%% The fourth publication (yet another journal paper, submitted for publication, note the optional parameter)
%% Note that you are allowed to use this option only when submitting the dissertation for pre-examination!
% Set the publication information, detailed information is not printed
\addpublication[submitted]{Journal Paper 3 Authors}{Journal Paper 3 Title}{Journal Name 3}{}{Submission date}{Year}{No copyright holder at this moment}{j3}
% Add the dissertation author's contribution to that publication.
\addcontribution{The author did everything}
% Add the errata of the publication, remove if there are none. (in submitted paper this is unlikely)
\adderrata{j3 IV This is also very wrong}
% Add the publication pdf file, the filename is the parameter.
\addpublicationpdf{dummyarticles/dummypdfarticle3.pdf}
\end{document}